\chapter{Glossar}\label{ch:glossar}

\section*{A}

\begin{description}
    \item[AABB (Axis-Aligned Bounding Box)] Achsenausgerichtete Begrenzungsbox, die ein Objekt über Minima und Maxima entlang der Raumachsen beschreibt.
    In \textit{helios} wird die \textit{AABB} unter anderem für effiziente Kollisionsprüfungen verwendet.
\end{description}

\section*{C}

\begin{description}
    \item[CommandBuffer] Zentraler Nachrichtenpuffer, in den Systeme und Manager Commands einstellen.
    Die Verarbeitung erfolgt an definierten Synchronisationspunkten durch die jeweils zuständigen \texttt{CommandHandler}.

    \item[CommandHandler] Handler-Konzept zur Verarbeitung bestimmter Command-Typen.
    Die Registrierung erfolgt in \textit{helios} typbasiert und nicht über Vererbung, \textit{Concepts} werden zur Kompilierzeit zur Validierung der Handler genutzt.

    \item[CommitPoint] Synchronisationspunkt innerhalb einer \textit{GameLoop}-Phase.
    An einem \texttt{CommitPoint} werden Buffer-Swaps durchgeführt, sodass nachfolgende Passes einen wohldefinierten Nachrichtenstand lesen können.

    \item[Component] Datenhaltender Baustein eines \texttt{GameObject}s bzw. einer Entität.
    Komponenten beschreiben in \textit{helios} fachliche Aspekte wie Position, Zustand oder Kollisionsform und werden in \textit{Sparse Sets} gespeichert.

    \item[Concept] C++ Sprachfeature, mit dem sich Anforderungen an Typen in Template-Kontexten deklarativ formulieren und zur Kompilierzeit prüfen lassen.
    In \textit{helios} werden \textit{Concepts} unter anderem dazu verwendet, um Rollen und Schnittstellen von Laufzeitkomponenten zu validieren.
\end{description}

\section*{E}

\begin{description}
    \item[Entität (Entity)] Abstrakter Begriff für Objekte, die in der GameWorld zur Laufzeit existieren. Auch als \textit{GameObject} bezeichnet.

    \item[Entity Component System (ECS)] Architekturmuster, das Objekte nicht über Klassenhierarchien, sondern über die Zusammensetzung aus Komponenten und darauf operierenden Systemen beschreibt.
    In \textit{helios} bildet ECS das zentrale Ordnungsprinzip der Laufzeitarchitektur.

    \item[EntityHandle] Zusammengesetzter Schlüssel aus \texttt{EntityId} und \texttt{VersionId}.
    Er dient in \textit{helios} zur eindeutigen Identifizierung von Entitäten und ermöglicht \textit{Generational Indexing}.

    \item[EntityManager] Klasse zur Verwaltung von Entitäten und ihren komponentenspezifischen Speicherstrukturen.
    Der EntityManager koordiniert unter anderem die \textit{Sparse Sets} und den Lebenszyklus von \texttt{EntityHandle}s.

    \item[Event/EventBus] Nachrichten, die zwischen Systemen und Laufzeitkomponenten ausgetauscht werden.
    In \textit{helios} wird zwischen Pass-, Phase- und Frame-Events unterschieden.

    \item[EntityRegistry] Datenstruktur zur Lifecycle-Verwaltung von Entitäten.
    Die EntityRegistry hält in \textit{helios} die Informationen zu \texttt{EntityHandle}s, die in er GameWorld existieren.
    Hierzu gehören Versionsstände und freigegebene Indizes.
    Auf Anfrage erzeugt sie neue EntityHandles, dabei können auch freigewordene \textit{EntityId}s wiederverwendet werden.
\end{description}

\section*{G}

\begin{description}
    \item[GameLoop] Zeitliche Ordnungsinstanz der Engine.
    In \textit{helios} ist die GameLoop in \textit{Phasen} und \textit{Passes} gegliedert und ermöglicht so die deterministische Ausführung von Systemen sowie den Austausch von Nachrichten über Systemgrenzen hinweg.

    \item[GameObject] Leichtgewichtige Fassade für den einfachen Zugriff auf ein in der GameWorld existierendes Entität.
    In \textit{helios} kapselt ein \texttt{GameObject} im Wesentlichen ein \texttt{EntityHandle} sowie den Zugriff auf den \texttt{EntityManager} und bietet eine semantisch klare API für die Abfrage und Manipulation von Entitäten an.

    \item[GameState] Übergeordneter Zustandsautomat der Anwendung.
    Der GameState beschreibt globale Zustände wie \texttt{Title}, \texttt{Running} oder \texttt{Paused}.

    \item[GameWorld] Zentraler Container zur Verwaltung von Ressourcen udn Laufzeitkompoenenten.
    Die GameWorld kapselt den Zugriff auf den \texttt{EntityManager}, besitzt die Registries für alle Manager und CommandHandler und bietet Zugriff auf Views und Session-Daten an.

    \item[Generational Indexing] Technik zur sicheren Wiederverwendung von Entity-Indizes.
    Durch die Kombination aus Entitätsindex und Versionsindex können veraltete EntityHandles (\textit{stale handles}) erkannt und bei Bedarf wiederverwendet werden.
    Bildet die zentrale Grundlage zur Realisierung von dichten Speicherstrukturen (\textit{dense layout}).
\end{description}

\section*{H}
    \begin{description}
    \item[Hook] Erweiterungspunkt (\textit{Extension Point}) für den Aufruf innerhalb eines fest definierten Ablaufs.
    Komponenten über Hooks wie \texttt{onAcquire}, \texttt{onRelease}, \texttt{onActivate} oder \texttt{onClone} auf bestimmte Zustandswechsel reagieren, ohne dass dafür eine gemeinsame Basisklasse erforderlich ist.
    Wird in \textit{helios} auch als allgemeiner Begriff für die Methoden eines Listeners genutzt.
\end{description}

\section*{M}

\begin{description}
    \item[Manager] Laufzeitkomponente zur Orchestrierung bestimmter Aufgabenbereiche.
    Manager fungieren in \textit{helios} oft auch als \texttt{CommandHandler} und führen gesammelte Befehle zu definierten Zeitpunkten aus.

    \item[MatchState] Unterhalb des \texttt{GameState} angesiedelter Zustandsautomat, der den Ablauf eines konkreten Spieldurchlaufs beschreibt, etwa über Zustände wie \texttt{Warmup}, \texttt{Playing} oder \texttt{GameOver}.
\end{description}

\section*{P}

\begin{description}
    \item[Pass] Logische Untereinheit einer \textit{GameLoop}-Phase, in der ein bestimmter Satz von Systemen ausgeführt wird.

    \item[Phase] Übergeordnete Gliederung der \textit{GameLoop}.
    In \textit{helios} wird zwischen \texttt{Pre}, \texttt{Main} und \texttt{Post} unterschieden.

    \item[Pooling / Object Pool] Technik zur Wiederverwendung bereits erzeugter Objekte.
    In \textit{helios} wird Pooling eingesetzt, um Speicherallokationen zur Laufzeit zu reduzieren und die Wiederverwendung von \texttt{GameObject}s zu ermöglichen, die über ein gemeinsames Prefab erstellt wurden.

    \item[Prefab] Blaupause für ein \texttt{GameObject}.
    Ein Prefab beschreibt eine Objektvorlage aus einer Zusammensetzung verschiedener Komponenten.
    In \textit{helios} werden Prefabs über Builder- bzw. Factory-Mechanismen definiert und anschließend für Spawn- und Pooling-Prozesse wiederverwendet.
\end{description}

\section*{S}

\begin{description}
    \item[Session] Spezielles \texttt{GameObject}, das Komponenten für globale laufzeitbezogene und spielspezifische Informationen vorhält, wie Zustände oder aktive Viewports.

    \item[SparseSet] Zentrale Datenstruktur der ECS-Implementierung in \textit{helios}.
    Sie bildet einen spärlich besetzten (\textit{sparse}) diskreten Wertebereich auf ein dichtes Array ab und ermöglicht so effizienten Zugriff, Einfügen und Entfernen von Komponenten in $\mathcal{O}(1)$.

    \item[StateManager] Manager zur Verwaltung regelbasierter Zustandsautomaten.
    Ein StateManager verarbeitet Zustandswechsel und benachrichtigt registrierte Listener über entsprechende Hooks.

    \item[Structure-of-Arrays (SoA)] Speicherlayout, bei dem gleichartige Daten in jeweils eigenen, dichten Arrays abgelegt werden.
    Dieses Layout bildet in \textit{helios} die Grundlage der komponentenorientierten Datenhaltung.

    \item[System] Verhaltensträger im ECS.
    Systeme iterieren über Mengen von Entitäten mit bestimmten Komponenten und aktualisieren deren Zustand.
\end{description}

\section*{T}

\begin{description}
    \item[Tag Component] Komponente ohne fachliche Nutzdaten, die ausschließlich der Markierung dient, etwa \texttt{Active} oder \texttt{DeadTagComponent}.

    \item[TypeIndex] Dichter Index zur Zuordnung einer Position zu einem Typen innerhalb einer vordefinierten Typmenge.
    Der TypeIndex dient in \textit{helios} der effizienten Identifizierung von Typen genutzt und ermöglicht so einen Zugriff in $\mathcal{O}(1)$.
\end{description}

\section*{U}

\begin{description}
    \item[UpdateContext] Kontextobjekt, das Systemen während der Ausführung in der \textit{GameLoop} übergeben wird.
    Der UpdateContext wird in der GameLoop als Parameter in tiefere Schichten weitergegeben und vereint den Zugriff auf Ressourcen und Laufzeitkomponenten, Event-Busse, Session-Daten, Timing-Informationen und Views.
\end{description}

\section*{V}

\begin{description}
    \item[View] Sicht auf Entitäten mit einer bestimmten Komponenten-Konstellation.
    Views werden in \textit{helios} von Systemen verwendet, um relevante \texttt{GameObject}s effizient zu iterieren.
\end{description}